\documentclass[10pt,a4paper]{amsart}
\usepackage{amsmath,amssymb,amsthm,graphicx}
% Retained minimal preamble for the Stacks Project chapter varieties.tex.
% Verbatim excerpt of the chapter preamble: Xy-pic setup (native preamble lines 14-18), the native theorem
% declarations with the mathematical macros (native preamble lines 154-217), and the project own no-op
% definitions of the reference/slogan/history display environments (native wrapper.tex lines 1-3) that the
% counted statement uses for its slogan. Copied mechanically; nothing added, removed or redefined.
% For commutative diagrams we use Xy-pic
\usepackage[all]{xy}

% We use 2cell for 2-commutative diagrams.
\xyoption{2cell}
% Theorem environments.
%
\theoremstyle{plain}
\newtheorem{theorem}[subsection]{Theorem}
\newtheorem{proposition}[subsection]{Proposition}
\newtheorem{lemma}[subsection]{Lemma}

\theoremstyle{definition}
\newtheorem{definition}[subsection]{Definition}
\newtheorem{example}[subsection]{Example}
\newtheorem{exercise}[subsection]{Exercise}
\newtheorem{situation}[subsection]{Situation}

\theoremstyle{remark}
\newtheorem{remark}[subsection]{Remark}
\newtheorem{remarks}[subsection]{Remarks}

\numberwithin{equation}{subsection}

% Macros
%
\def\lim{\mathop{\mathrm{lim}}\nolimits}
\def\colim{\mathop{\mathrm{colim}}\nolimits}
\def\Spec{\mathop{\mathrm{Spec}}}
\def\Hom{\mathop{\mathrm{Hom}}\nolimits}
\def\Ext{\mathop{\mathrm{Ext}}\nolimits}
\def\SheafHom{\mathop{\mathcal{H}\!\mathit{om}}\nolimits}
\def\SheafExt{\mathop{\mathcal{E}\!\mathit{xt}}\nolimits}
\def\Sch{\mathit{Sch}}
\def\Mor{\mathop{\mathrm{Mor}}\nolimits}
\def\Ob{\mathop{\mathrm{Ob}}\nolimits}
\def\Sh{\mathop{\mathit{Sh}}\nolimits}
\def\NL{\mathop{N\!L}\nolimits}
\def\CH{\mathop{\mathrm{CH}}\nolimits}
\def\proetale{{pro\text{-}\acute{e}tale}}
\def\etale{{\acute{e}tale}}
\def\QCoh{\mathit{QCoh}}
\def\Ker{\mathop{\mathrm{Ker}}}
\def\Im{\mathop{\mathrm{Im}}}
\def\Coker{\mathop{\mathrm{Coker}}}
\def\Coim{\mathop{\mathrm{Coim}}}

% Boxtimes
%
\DeclareMathSymbol{\boxtimes}{\mathbin}{AMSa}{"02}

%
% Macros for moduli stacks/spaces
%
\def\QCohstack{\mathcal{QC}\!\mathit{oh}}
\def\Cohstack{\mathcal{C}\!\mathit{oh}}
\def\Spacesstack{\mathcal{S}\!\mathit{paces}}
\def\Quotfunctor{\mathrm{Quot}}
\def\Hilbfunctor{\mathrm{Hilb}}
\def\Curvesstack{\mathcal{C}\!\mathit{urves}}
\def\Polarizedstack{\mathcal{P}\!\mathit{olarized}}
\def\Complexesstack{\mathcal{C}\!\mathit{omplexes}}
% \Pic is the operator that assigns to X its picard group, usage \Pic(X)
% \Picardstack_{X/B} denotes the Picard stack of X over B
% \Picardfunctor_{X/B} denotes the Picard functor of X over B
\def\Pic{\mathop{\mathrm{Pic}}\nolimits}
\def\Picardstack{\mathcal{P}\!\mathit{ic}}
\def\Picardfunctor{\mathrm{Pic}}
\def\Deformationcategory{\mathcal{D}\!\mathit{ef}}
\newenvironment{reference}{}{}
\newenvironment{slogan}{}{}
\newenvironment{history}{}{}
\usepackage{cleveref}
\crefname{originalreference}{original reference}{original references}
\Crefname{originalreference}{Original reference}{Original references}
\makeatletter
\newlabel{morphisms-lemma-finite-type-over-affine-ample-very-ample}{{01VT}{source}{}{sourcepointer-1}{}}
\newlabel{morphisms-lemma-finite-type-over-affine-ample-very-ample@cref}{{[originalreference][0][]01VT}{[1][0][]source}{}{}{}}
\makeatother
\crefname{originalreference}{original reference}{original references}
\Crefname{originalreference}{Original reference}{Original references}
\makeatletter
\newlabel{coherent-lemma-relative-affine-cohomology}{{089W}{source}{}{sourcepointer-2}{}}
\newlabel{coherent-lemma-relative-affine-cohomology@cref}{{[originalreference][0][]089W}{[1][0][]source}{}{}{}}
\makeatother
\crefname{originalreference}{original reference}{original references}
\Crefname{originalreference}{Original reference}{Original references}
\makeatletter
\newlabel{algebra-lemma-depth-goes-down-finite}{{0AUK}{source}{}{sourcepointer-3}{}}
\newlabel{algebra-lemma-depth-goes-down-finite@cref}{{[originalreference][0][]0AUK}{[1][0][]source}{}{}{}}
\makeatother
\crefname{originallemma}{original lemma}{original lemmas}
\Crefname{originallemma}{Original lemma}{Original lemmas}
\makeatletter
\newlabel{lemma-exact-sequence-induction}{{08A0}{source}{}{sourcepointer-4}{}}
\newlabel{lemma-exact-sequence-induction@cref}{{[originallemma][0][]08A0}{[1][0][]source}{}{}{}}
\makeatother
\crefname{originalreference}{original reference}{original references}
\Crefname{originalreference}{Original reference}{Original references}
\makeatletter
\newlabel{algebra-lemma-depth-drops-by-one}{{090R}{source}{}{sourcepointer-5}{}}
\newlabel{algebra-lemma-depth-drops-by-one@cref}{{[originalreference][0][]090R}{[1][0][]source}{}{}{}}
\makeatother
\crefname{originallemma}{original lemma}{original lemmas}
\Crefname{originallemma}{Original lemma}{Original lemmas}
\makeatletter
\newlabel{lemma-vanishin-h0-negative}{{0FD7}{source}{}{sourcepointer-6}{}}
\newlabel{lemma-vanishin-h0-negative@cref}{{[originallemma][0][]0FD7}{[1][0][]source}{}{}{}}
\makeatother
\crefname{originalremark}{original remark}{original remarks}
\Crefname{originalremark}{Original remark}{Original remarks}
\makeatletter
\newlabel{remark-exact-sequence-induction-cohomology}{{0EGK}{source}{}{sourcepointer-7}{}}
\newlabel{remark-exact-sequence-induction-cohomology@cref}{{[originalremark][0][]0EGK}{[1][0][]source}{}{}{}}
\makeatother
\crefname{originalreference}{original reference}{original references}
\Crefname{originalreference}{Original reference}{Original references}
\makeatletter
\newlabel{coherent-lemma-flat-base-change-cohomology}{{02KH}{source}{}{sourcepointer-8}{}}
\newlabel{coherent-lemma-flat-base-change-cohomology@cref}{{[originalreference][0][]02KH}{[1][0][]source}{}{}{}}
\makeatother
\crefname{originallemma}{original lemma}{original lemmas}
\Crefname{originallemma}{Original lemma}{Original lemmas}
\makeatletter
\newlabel{lemma-change-fields-flat}{{0C4Y}{source}{}{sourcepointer-9}{}}
\newlabel{lemma-change-fields-flat@cref}{{[originallemma][0][]0C4Y}{[1][0][]source}{}{}{}}
\makeatother
\crefname{originalreference}{original reference}{original references}
\Crefname{originalreference}{Original reference}{Original references}
\makeatletter
\newlabel{algebra-lemma-apply-grothendieck-module}{{0338}{source}{}{sourcepointer-10}{}}
\newlabel{algebra-lemma-apply-grothendieck-module@cref}{{[originalreference][0][]0338}{[1][0][]source}{}{}{}}
\makeatother
\title{Enriques-Severi-Zariski: vanishing of H\textasciicircum{}1 of sufficiently negative twists of an ample invertible sheaf under depth at least two}
\author{The Stacks Project Authors}
\date{stacks-project:varieties.tex:local-snapshot-20261011; GNU FDL 1.2 or later}
\begin{document}
\maketitle
\noindent\textit{Editorial note.} All delivered mathematical text is the author's own, copied verbatim from the pinned Stacks Project chapter file varieties.tex; no statement or proof step was rewritten, reordered or removed.\par\medskip
\noindent\textit{Editorial note.} The counted span cites ten cross-chapter labels, each recorded in assembly.external\_references with its canonical tag as a navigation pointer.\par\medskip
\noindent\textit{Editorial note.} The proof cites no bibliography item inside the counted span; the author's attribution to Enriques-Severi-Zariski is part of the lemma's own name and was copied verbatim.\par\medskip
\noindent\textit{Allowed standard prerequisites (editorial).}\begin{itemize}
\item Very ample invertible sheaves and closed immersions into projective space. Stacks Project, Morphisms, Lemma morphisms-lemma-finite-type-over-affine-ample-very-ample (tag 01VR). Produces the embedding making L a power of the restriction of O(1).
\item Cohomology of a closed-immersion pushforward and depth along finite ring maps. Stacks Project, Cohomology of Schemes, Lemma coherent-lemma-relative-affine-cohomology (tag 089W); Algebra, Lemma algebra-lemma-depth-goes-down-finite (tag 0AUH). Transfer cohomology and depth between F and the packaged sheaf G.
\end{itemize}\medskip
\ifdefined\hypertarget\hypertarget{sourcepointer-1}{}\fi
\noindent\textit{Source reference.} Original reference 01VT: sources/stacks-chow/tags.txt, source lines 2458--2458. Cross-chapter standard prerequisite cited by the author inside the counted span; it is used as a pointer and its own proof is not part of this item.\par\medskip
\ifdefined\hypertarget\hypertarget{sourcepointer-2}{}\fi
\noindent\textit{Source reference.} Original reference 089W: sources/stacks-chow/tags.txt, source lines 10511--10511. Cross-chapter standard prerequisite cited by the author inside the counted span; it is used as a pointer and its own proof is not part of this item.\par\medskip
\ifdefined\hypertarget\hypertarget{sourcepointer-3}{}\fi
\noindent\textit{Source reference.} Original reference 0AUK: sources/stacks-chow/tags.txt, source lines 13653--13653. Cross-chapter standard prerequisite cited by the author inside the counted span; it is used as a pointer and its own proof is not part of this item.\par\medskip
\ifdefined\hypertarget\hypertarget{sourcepointer-4}{}\fi
\noindent\textit{Source reference.} Original lemma 08A0: sources/stacks-chow/tags.txt, source lines 10515--10515. Cross-chapter standard prerequisite cited by the author inside the counted span; it is used as a pointer and its own proof is not part of this item.\par\medskip
\ifdefined\hypertarget\hypertarget{sourcepointer-5}{}\fi
\noindent\textit{Source reference.} Original reference 090R: sources/stacks-chow/tags.txt, source lines 11419--11419. Cross-chapter standard prerequisite cited by the author inside the counted span; it is used as a pointer and its own proof is not part of this item.\par\medskip
\ifdefined\hypertarget\hypertarget{sourcepointer-6}{}\fi
\noindent\textit{Source reference.} Original lemma 0FD7: sources/stacks-chow/tags.txt, source lines 19207--19207. Cross-chapter standard prerequisite cited by the author inside the counted span; it is used as a pointer and its own proof is not part of this item.\par\medskip
\ifdefined\hypertarget\hypertarget{sourcepointer-7}{}\fi
\noindent\textit{Source reference.} Original remark 0EGK: sources/stacks-chow/tags.txt, source lines 18100--18100. Cross-chapter standard prerequisite cited by the author inside the counted span; it is used as a pointer and its own proof is not part of this item.\par\medskip
\ifdefined\hypertarget\hypertarget{sourcepointer-8}{}\fi
\noindent\textit{Source reference.} Original reference 02KH: sources/stacks-chow/tags.txt, source lines 3372--3372. Cross-chapter standard prerequisite cited by the author inside the counted span; it is used as a pointer and its own proof is not part of this item.\par\medskip
\ifdefined\hypertarget\hypertarget{sourcepointer-9}{}\fi
\noindent\textit{Source reference.} Original lemma 0C4Y: sources/stacks-chow/tags.txt, source lines 15243--15243. Cross-chapter standard prerequisite cited by the author inside the counted span; it is used as a pointer and its own proof is not part of this item.\par\medskip
\ifdefined\hypertarget\hypertarget{sourcepointer-10}{}\fi
\noindent\textit{Source reference.} Original reference 0338: sources/stacks-chow/tags.txt, source lines 4047--4047. Cross-chapter standard prerequisite cited by the author inside the counted span; it is used as a pointer and its own proof is not part of this item.\par\medskip
\begin{lemma}[Enriques-Severi-Zariski]
\label{lemma-vanishin-h1-negative}
Let $k$ be a field. Let $X$ be a proper scheme over $k$. Let $\mathcal{L}$
be an ample invertible $\mathcal{O}_X$-module. Let $\mathcal{F}$ be a
coherent $\mathcal{O}_X$-module. Assume that
for $x \in X$ closed we have $\text{depth}(\mathcal{F}_x) \geq 2$.
Then
$H^1(X, \mathcal{F} \otimes_{\mathcal{O}_X} \mathcal{L}^{\otimes m}) = 0$
for $m \ll 0$.
\end{lemma}

\begin{proof}
Choose a closed immersion $i : X \to \mathbf{P}^n_k$ such that
$i^*\mathcal{O}(1) \cong \mathcal{L}^{\otimes e}$ for some $e > 0$
(see Morphisms, Lemma
\ref{morphisms-lemma-finite-type-over-affine-ample-very-ample}).
Then it suffices to prove the lemma for
$$
\mathcal{G} =
i_*(\mathcal{F} \oplus \mathcal{F} \otimes \mathcal{L} \oplus \ldots
\oplus \mathcal{F} \otimes \mathcal{L}^{\otimes e - 1})
\quad\text{and}\quad \mathcal{O}(1)
$$
on $\mathbf{P}^n_k$. Namely, we have
$$
H^1(\mathbf{P}^n_k, \mathcal{G}(m)) =
\bigoplus\nolimits_{j = 0, \ldots, e - 1}
H^1(X, \mathcal{F} \otimes \mathcal{L}^{\otimes j + me})
$$
by
Cohomology of Schemes, Lemma \ref{coherent-lemma-relative-affine-cohomology}.
Also, if $y \in \mathbf{P}^n_k$ is a closed point then
$\text{depth}(\mathcal{G}_y) = \infty$ if $y \not \in i(X)$
and $\text{depth}(\mathcal{G}_y) = \text{depth}(\mathcal{F}_x)$
if $y = i(x)$ because in this case
$\mathcal{G}_y \cong \mathcal{F}_x^{\oplus e}$ as a module over
$\mathcal{O}_{\mathbf{P}^n_k, x}$ and we can use for example
Algebra, Lemma \ref{algebra-lemma-depth-goes-down-finite}
to get the equality.

\medskip\noindent
Assume $X = \mathbf{P}^n_k$ and $\mathcal{L} = \mathcal{O}(1)$ and
$k$ is infinite. Choose $s \in H^0(\mathbf{P}^1_k, \mathcal{O}(1))$
which determines an exact sequence
$$
0 \to \mathcal{F}(-1) \xrightarrow{s} \mathcal{F} \to \mathcal{G} \to 0
$$
as in Lemma \ref{lemma-exact-sequence-induction}. Since the map
$\mathcal{F}(-1) \to \mathcal{F}$ is affine locally given by
multiplying by a nonzerodivisor on $\mathcal{F}$
we see that for $x \in \mathbf{P}^n_k$ closed we have
$\text{depth}(\mathcal{G}_x) \geq 1$, see
Algebra, Lemma \ref{algebra-lemma-depth-drops-by-one}.
Hence by Lemma \ref{lemma-vanishin-h0-negative}
we have $H^0(\mathcal{G}(m)) = 0$ for $m \ll 0$.
Looking at the long exact sequence of cohomology after twisting
(see Remark \ref{remark-exact-sequence-induction-cohomology})
we find that the sequence of numbers
$$
\dim H^1(\mathbf{P}^n_k, \mathcal{F}(m))
$$
stabilizes for $m \leq m_0$ for some integer $m_0$.
Let $N$ be the common dimension of these
spaces for $m \leq m_0$. We have to show $N = 0$.

\medskip\noindent
For $d > 0$ and $m \leq m_0$ consider the bilinear map
$$
H^0(\mathbf{P}^n_k, \mathcal{O}(d)) \times
H^1(\mathbf{P}^n_k, \mathcal{F}(m - d))
\longrightarrow
H^1(\mathbf{P}^n_k, \mathcal{F}(m))
$$
By linear algebra, there is a codimension $\leq N^2$ subspace
$V_m \subset H^0(\mathbf{P}^n_k, \mathcal{O}(d))$ such that
multiplication by $s' \in V_m$ annihilates
$H^1(\mathbf{P}^n_k, \mathcal{F}(m - d))$.
Observe that for $m' < m \leq m_0$ the diagram
$$
\xymatrix{
H^0(\mathbf{P}^n_k, \mathcal{O}(d)) \times
H^1(\mathbf{P}^n_k, \mathcal{F}(m' - d)) \ar[r] \ar[d]^{1 \times s^{m' - m}} &
H^1(\mathbf{P}^n_k, \mathcal{F}(m')) \ar[d]^{s^{m' - m}}\\
H^0(\mathbf{P}^n_k, \mathcal{O}(d)) \times
H^1(\mathbf{P}^n_k, \mathcal{F}(m - d)) \ar[r] &
H^1(\mathbf{P}^n_k, \mathcal{F}(m))
}
$$
commutes with isomorphisms going vertically. Thus $V_m = V$ is
independent of $m \leq m_0$. For $x \in \text{Ass}(\mathcal{F})$
set $Z = \overline{\{x\}}$. For $d$ large enough the linear map
$$
H^0(\mathbf{P}^n_k, \mathcal{O}(d)) \to H^0(Z, \mathcal{O}(d)|_Z)
$$
has rank $> N^2$ because $\dim(Z) \geq 1$ (for example this follows
from asymptotic Riemann-Roch and ampleness $\mathcal{O}(1)$; details
omitted). Hence we can find $s' \in V$ such that $s'$ does not vanish
in any associated point of $\mathcal{F}$ (use that the set
of associated points is finite). Then we obtain
$$
0 \to \mathcal{F}(-d) \xrightarrow{s'} \mathcal{F} \to \mathcal{G}' \to 0
$$
and as before we conclude as before that multiplication by $s'$
on $H^1(\mathbf{P}^n_k, \mathcal{F}(m - d))$ is injective
for $m \ll 0$. This contradicts the choice of $s'$ unless
$N = 0$ as desired.

\medskip\noindent
We still have to treat the case where $k$ is finite.
In this case let $K/k$ be any infinite algebraic field extension.
Denote $\mathcal{F}_K$ and $\mathcal{L}_K$ the pullbacks
of $\mathcal{F}$ and $\mathcal{L}$ to $X_K = \Spec(K) \times_{\Spec(k)} X$.
We have
$$
H^1(X_K, \mathcal{F}_K \otimes \mathcal{L}_K^{\otimes m}) =
H^1(X, \mathcal{F} \otimes \mathcal{L}^{\otimes m}) \otimes_k K
$$
by Cohomology of Schemes, Lemma
\ref{coherent-lemma-flat-base-change-cohomology}.
On the other hand, a closed point $x_K$ of $X_K$ maps to a closed point
$x$ of $X$ because $K/k$ is an algebraic extension. The
ring map $\mathcal{O}_{X, x} \to \mathcal{O}_{X_K, x_K}$
is flat (Lemma \ref{lemma-change-fields-flat}). Hence we have
$$
\text{depth}(\mathcal{F}_{x_K}) =
\text{depth}(\mathcal{F}_x \otimes_{\mathcal{O}_{X, x}}
\mathcal{O}_{X_K, x_K}) \geq
\text{depth}(\mathcal{F}_x)
$$
by Algebra, Lemma \ref{algebra-lemma-apply-grothendieck-module}
(in fact equality holds here but we don't need it).
Therefore the result over $k$ follows
from the result over the infinite field $K$ and the proof is complete.
\end{proof}

\end{document}
